\section{逆矩阵求导应写成线性求解}
\label{sec:10-Matrix-Calculus/03-Trace-Calculus/02}

一个隐式层要接进网络：给定参数矩阵 \(A\) 和输入 \(b\)，这一层输出线性系统 \(Ax=b\) 的解，下游再接损失 \(\ell(x)\)。前向部分没什么可说的，调一次分解就完事。反向卡住了：需要 \(\nabla_A\ell\)，而 \(A\) 有 \(n^2\) 个元素。按公式把 \(x\) 写成 \(A^{-1}b\)、对 \(A^{-1}\) 求导、再把结果乘出来，代码能跑，但 \(n=4000\) 时反向耗时是前向的几十倍，显存也顶到上限。

问题不在公式。公式只有一行，而且不需要背。只要 \(X\) 可逆，恒等式 \(XX^{-1}=I\) 两边微分给出 \((dX)X^{-1}+X\,d(X^{-1})=0\)，左乘 \(X^{-1}\) 得

\[
d(X^{-1})=-X^{-1}(dX)X^{-1}.
\]

负号来自移项，两个 \(X^{-1}\) 一个来自消去左侧的 \(X\)、一个是原式里本来就有的。记不清的时候重新微分恒等式，比回忆符号可靠。

麻烦出在把这一行翻译成代码的方式。式子里写着三个 \(X^{-1}\)，而正确的实现里一个显式的逆都不出现。

\subsection{符号里的逆是记号，落到代码里要变成求解}

\(X^{-1}M\) 这个记号的含义是``\(XY=M\) 的解 \(Y\)''，它并不指示先造出 \(X^{-1}\) 再做一次矩阵乘法。数值分析那边已经把账算过：显式求逆要 \(O(n^3)\) 次运算、\(n^2\) 的存储，误差比直接求解大，而且丢掉了 \(X\) 可能有的稀疏或带状结构，见\hyperref[sec:08-Numerical-Analysis/07-Numerical-Methods-for-ML/03]{``需要的是解，而不是逆矩阵''}与\hyperref[sec:08-Numerical-Analysis/05-Numerical-Linear-Algebra/01]{``线性方程为何分解而不求逆''}。

这条纪律在求导场景里比在前向场景里更值钱，因为微分公式会把 \(X^{-1}\) 复制好几份。照着符号写代码的人会造出两三个逆；照着含义写代码的人，把同一份分解用两三次。

\subsection{对方程微分，而不是对解的表达式微分}

前向要算的是 \(x=A^{-1}b\)，也就是 \(Ax=b\) 的解。求导时不必先把解写成含逆的表达式，直接对方程两边微分：

\[
(dA)x+A\,dx=db
\qquad\Longrightarrow\qquad
A\,dx=db-(dA)x.
\]

给定扰动 \(dA\) 与 \(db\)，右端是一个能直接算出来的向量，\(dx\) 就是以 \(A\) 为系数矩阵的又一个线性系统的解。系数矩阵和前向那次完全相同，如果前向已经做了 \(LU\) 或 Cholesky 分解，这里只需一次回代。

这种写法叫隐式微分。它比``代入 \(x=A^{-1}b\) 再套逆矩阵公式''少绕一圈，而且从形式上就杜绝了显式求逆——推导过程中根本没有 \(A^{-1}\) 这个符号出现的机会。

\subsection{反向只需一次转置求解和一个外积}

上面那条路径是前向模式：给一个输入扰动，问输出怎么变。训练要的是反向模式：给定上游梯度 \(\bar x=\nabla_x\ell\)，问参数该怎么改。把 \(dx=A^{-1}\bigl(db-(dA)x\bigr)\) 代进 \(d\ell=\bar x^{\trans}dx\)：

\[
d\ell=\bar x^{\trans}A^{-1}db-\bar x^{\trans}A^{-1}(dA)x.
\]

两项里都出现同一个组合 \(\bar x^{\trans}A^{-1}\)。给它一个名字，令 \(\lambda\) 满足

\[
A^{\trans}\lambda=\bar x,
\]

于是 \(\lambda^{\trans}=\bar x^{\trans}A^{-1}\)，微分变成 \(d\ell=\lambda^{\trans}db-\lambda^{\trans}(dA)x\)。第二项用上一节的读法整理：\(\lambda^{\trans}(dA)x=\trace\bigl((\lambda x^{\trans})^{\trans}dA\bigr)\)。两个梯度随即读出：

\[
\nabla_b\ell=\lambda,
\qquad
\nabla_A\ell=-\lambda x^{\trans}.
\]

\(\lambda\) 被称为伴随变量。整个反向过程只做了一件计算量可观的事：解一次转置系统。其余是一个外积。

\begin{center}
\begin{tikzpicture}[font=\small,>=stealth]
  \node[font=\footnotesize] at (-0.55,1.15) {\(b\)};
  \node[draw,black!55,rounded corners,fill=blue!4,minimum height=8mm,inner sep=5pt] (F) at (2.1,1.15) {解 \(Ax=b\)};
  \node[font=\footnotesize] at (4.3,1.15) {\(x\)};
  \node[draw,black!55,rounded corners,minimum height=8mm,inner sep=5pt] (L) at (6.2,1.15) {\(\ell(x)\)};
  \draw[->,black!65] (-0.25,1.15) -- (F);
  \draw[->,black!65] (F) -- (4.05,1.15);
  \draw[->,black!65] (4.55,1.15) -- (L);
  \node[font=\footnotesize,black!60,anchor=west] at (6.85,1.15) {前向};

  \node[font=\footnotesize] at (-0.9,-1.15) {\(\nabla_A\ell=-\lambda x^{\trans}\)};
  \node[draw,red!60!black,rounded corners,fill=red!4,minimum height=8mm,inner sep=5pt] (Bk) at (2.1,-1.15) {解 \(A^{\trans}\lambda=\bar x\)};
  \node[font=\footnotesize] at (4.3,-1.15) {\(\bar x\)};
  \draw[<-,red!60!black] (0.35,-1.15) -- (Bk);
  \draw[<-,red!60!black] (Bk) -- (4.05,-1.15);
  \draw[<-,red!60!black] (4.55,-1.15) -- (L);
  \node[font=\footnotesize,red!60!black,anchor=west] at (6.85,-1.15) {反向};

  \draw[dashed,black!45,->] (F) -- (Bk);
  \node[font=\footnotesize,black!60,align=left,anchor=east] at (1.95,0) {复用同一份分解};
  \draw[dashed,black!45,->] (4.3,0.85) -- (4.3,-0.85);
  \node[font=\footnotesize,black!60,anchor=east] at (4.15,0) {前向的 \(x\) 要留着};
\end{tikzpicture}

\emph{前向解一次，反向解一次转置系统，两次共用一份分解；\(A^{-1}\) 从头到尾没有被构造出来。}
\end{center}

图里那条竖着的虚线是这一节的全部经济性所在。设 \(A\) 稠密，分解一次 \(\tfrac13 n^3\)，每次三角回代 \(n^2\)。前向付掉分解的钱之后，反向的边际成本是一次 \(n^2\) 的回代加一个 \(n^2\) 的外积——与 \(n^2\) 个参数的规模同阶，不可能更省。反过来，若照着符号显式构造 \(A^{-1}\)，光这一步就是另一个 \(n^3\)，而且它算出来的东西马上就要被乘掉、扔掉。

外积那一步还有一个实现上的提醒。\(\nabla_A\ell=-\lambda x^{\trans}\) 是稠密的 \(n\times n\) 矩阵，如果 \(A\) 本来是稀疏的、只有若干个非零位置由参数控制，就不该把这个外积整个物化，只取需要的那些位置即可。伴随方法在隐式层、Kalman 滤波、可微物理和 PDE 约束优化里都以这个形态出现，共同的模式是：前向解一个系统，反向解它的转置系统。

\subsection{含逆的二次型也走同一条路}

单独看一个纯符号推导。令 \(f(X)=a^{\trans}X^{-1}b\)，代入逆的微分公式，

\[
df=-a^{\trans}X^{-1}(dX)X^{-1}b.
\]

令 \(u=X^{-\trans}a\)、\(v=X^{-1}b\)，它就成了 \(df=-u^{\trans}(dX)v=-\trace\bigl((uv^{\trans})^{\trans}dX\bigr)\)，于是

\[
\nabla_Xf=-uv^{\trans}=-X^{-\trans}ab^{\trans}X^{-\trans}.
\]

最右边那个含逆的表达式适合印在纸上，中间那个 \(-uv^{\trans}\) 才是拿去实现的。\(u\) 和 \(v\) 分别是 \(X^{\trans}u=a\) 与 \(Xv=b\) 的解，两次回代复用同一份分解。这里也能看出秩的结构：梯度是秩一矩阵，无论 \(n\) 多大。

若 \(X\) 被约束为对称，梯度还要多一步。此时 \(dX\) 只能在对称矩阵里取值，\(\trace(G^{\trans}dX)\) 只感知 \(G\) 的对称部分，应取 \(\tfrac12(G+G^{\trans})\) 作为约束下的梯度——与上一节``二次型只感知 \(A+A^{\trans}\)''是同一件事。当 \(X\) 对称正定且 \(a=b\) 时，无约束梯度 \(-X^{-1}bb^{\trans}X^{-1}\) 本来就对称，这一步是恒等变形，但把它当成普遍规律会在别处出错。

\subsection{可逆、条件数与秩变化是三道前提}

这一节的每个公式都要求 \(A\) 在考察点附近可逆。恒等式 \(XX^{-1}=I\) 在 \(X\) 奇异处根本没有定义，微分自然也无从谈起。

接近奇异是另一回事：公式仍然成立，但数值上不可信。条件数 \(\kappa(A)=\norm{A}\norm{A^{-1}}\) 同时刻画前向解和梯度对扰动的敏感度，\(\kappa\) 大的时候两者都会被放大，见\hyperref[sec:08-Numerical-Analysis/01-Floating-Point-and-Error/03]{``条件数衡量问题本身的敏感性''}。这种放大属于问题本身，换一个自动微分实现不会让它消失；能做的是在 \(A\) 上加阻尼项、换更稳的分解，或者承认梯度在这一带只有几位有效数字。

秩会变化的情形要单独对待。伪逆 \(A^{\dagger}\) 的导数不由上面的公式给出，把 \(A^{-1}\) 机械换成 \(A^{\dagger}\) 是错的。在秩固定的流形上可以推导出对应公式，但跨越秩变化点时映射本身就不光滑，那里没有导数可求。

最后是分解方式与精度语义。前向用 Cholesky 就在反向也用 Cholesky（见\hyperref[sec:08-Numerical-Analysis/05-Numerical-Linear-Algebra/03]{``Cholesky：利用正定结构''}），前向用迭代法解到容差 \(\varepsilon\)，反向的伴随系统也只能解到相当的容差，梯度随之只近似满足隐式方程。迭代求解器的停止准则因此不是一个纯粹的前向参数，它同时决定了梯度的可信位数；梯度检查在这类层上通不过，第一个该怀疑的就是这个容差，而不是推导。

逆矩阵是 \(X\) 的一种非多项式的出现方式，微分之后仍然是有理表达式。下一节的对象走得更远：它对 \(X\) 取行列式再取对数，微分之后却收缩回一个 trace，而它的数值实现同样不能照着符号写。
